%==============================================================================
%  GENERAL INTRODUCTION  (~1.5 pages, unnumbered chapter)
%==============================================================================
\chapter*{General Introduction}
\addcontentsline{toc}{chapter}{General Introduction}
\markboth{General Introduction}{General Introduction}

Machine learning has moved from research notebooks into the core of modern
software products. Yet the gap between \emph{training a model that works on a
laptop} and \emph{running that model reliably in production} remains one of the
hardest problems in the field. Studies have repeatedly shown that only a small
fraction of the effort in a real machine-learning system goes into the model
itself; the overwhelming majority is spent on data plumbing, configuration,
reproducibility, deployment, and monitoring~\cite{sculley2015hidden}. The
discipline that addresses this gap is \gls{mlops}, the application of
\gls{devops} principles to the machine-learning lifecycle.

In most teams, the path from experiment to deployed model is still manual and
fragile. Data scientists train models in notebooks, track results in
spreadsheets or not at all, and hand artifacts to engineers who deploy them
through ad-hoc scripts. The consequences are familiar: experiments that cannot
be reproduced, the recurring ``it works on my machine'' failure, and slow,
error-prone deployments. A unified platform that automates this lifecycle ---
training, experiment tracking, model registry, deployment, and monitoring ---
removes this friction and lets teams ship models with confidence.

This report presents the design and implementation of such a platform: an
end-to-end, self-hosted \gls{mlops} system carried out as an end-of-studies
project (PFE) within \textbf{INSOMEA}. The platform lets a user upload code and
data through a web interface, trigger a training pipeline that runs on a
\gls{k8s} cluster, automatically capture metrics and artifacts, register and
promote model versions, and deploy a chosen model as a live prediction service
, all from a single application, with continuous deployment to a public URL.

% ---- roadmap ---------------------------------------------------------------
The remainder of this report is organized as follows:

\begin{itemize}[leftmargin=1.4em]
  \item \textbf{Chapter~1, General Context} introduces the host organization,
        studies the existing situation and its limitations, states the problem,
        outlines the proposed solution, and presents the Scrum method and
        work environment adopted for the project.
  \item \textbf{Chapter~2, MLOps Background and Technology Foundations}
        defines the \gls{mlops} discipline, its core challenges, the enabling
        technologies used in this work, and positions the project against
        existing commercial solutions.
  \item \textbf{Chapter~3, Requirements Analysis and Project Planning}
        specifies the functional and non-functional requirements, presents the
        global use-case model, the product backlog, the release plan over four
        sprints, and the global system architecture.
  \item \textbf{Chapters~4 to~7} each document one sprint, following the same
        structure: sprint backlog, analysis and conception, realization, and
        sprint review.
  \item \textbf{General Conclusion} summarizes the results achieved, the skills
        gained, the limitations of the current system, and directions for future
        work.
\end{itemize}
